\subsection{Research design and paradigm}

\textbf{Design.}
- The study is a \textbf{controlled, between-architecture experiment with repeated randomised runs}: a full
factorial over connector × task with seeds as replicates, plus targeted follow-ups.
- Only one factor varies at a time: the vision-language connector, or one specified follow-up factor.
- Everything else is held fixed: the LLM, the frozen vision features, the data, the prompt, the optimiser and
the training budget.

\textbf{Why this design.}
- The research question is causal: \emph{does the connector change shortcut reliance?}
- Observational comparisons of released VLMs cannot answer it, because released models differ in data, scale,
encoder and training recipe at the same time.
- Design studies of VLMs \cite{mckinzie2024mm1,karamcheti2024prismatic} show that only controlled ablations can
attribute effects to one component.
- Seed-level replication is required because a single seed can reverse conclusions in deep learning
\cite{henderson2018matters,picard2021seed,bouthillier2021variance}.

\textbf{Paradigm.}
- The study is positivist and hypothesis-testing.
- The primary hypotheses (H1, H2), their tests, the significance level (one-sided α = 0.025) and the
equivalence margin (±3 WGA points) were written down before the seed-extension runs (see
\href{OUTLINE.md}{\texttt{OUTLINE.md}} and \href{../study/hypotheses.py}{\texttt{study/hypotheses.py}}).
- Analyses beyond these are labelled exploratory.

\subsection{Data sources, sampling and inclusion criteria}

{\def\LTcaptype{none} % do not increment counter
\begin{longtable}[]{@{}
  >{\raggedright\arraybackslash}p{(\linewidth - 12\tabcolsep) * \real{0.1290}}
  >{\raggedright\arraybackslash}p{(\linewidth - 12\tabcolsep) * \real{0.1290}}
  >{\raggedright\arraybackslash}p{(\linewidth - 12\tabcolsep) * \real{0.1290}}
  >{\raggedright\arraybackslash}p{(\linewidth - 12\tabcolsep) * \real{0.1290}}
  >{\centering\arraybackslash}p{(\linewidth - 12\tabcolsep) * \real{0.1613}}
  >{\centering\arraybackslash}p{(\linewidth - 12\tabcolsep) * \real{0.1613}}
  >{\centering\arraybackslash}p{(\linewidth - 12\tabcolsep) * \real{0.1613}}@{}}
\toprule\noalign{}
\begin{minipage}[b]{\linewidth}\raggedright
Task
\end{minipage} & \begin{minipage}[b]{\linewidth}\raggedright
Source
\end{minipage} & \begin{minipage}[b]{\linewidth}\raggedright
Shortcut (spurious attribute \emph{a})
\end{minipage} & \begin{minipage}[b]{\linewidth}\raggedright
Label \emph{y}
\end{minipage} & \begin{minipage}[b]{\linewidth}\centering
Train
\end{minipage} & \begin{minipage}[b]{\linewidth}\centering
Val
\end{minipage} & \begin{minipage}[b]{\linewidth}\centering
Test
\end{minipage} \\
\midrule\noalign{}
\endhead
\bottomrule\noalign{}
\endlastfoot
Waterbirds & CUB-200 birds on Places backgrounds \cite{wah2011cub,zhou2018places}, official splits \cite{sagawa2020groupdro} & background (land / water) & landbird / waterbird & 4,795 & 1,199 & 5,794 \\
Watermark (ours) & MS-COCO images \cite{lin2014coco} & stock-photo text watermark in the corner & dog present & 4,800 (95\% shortcut-aligned) & 1,200 (group-balanced) & 2,400 (group-balanced) \\
CelebA & CelebA, official partitions \cite{liu2015celeba} & gender (Male) & Blond\_Hair & 8,000 (stratified subsample) & 5,000 (≥ 100 per group) & full official test \\
VQA-CP v2 (open-ended) & VQA-CP v2 \cite{agrawal2018vqacp} on COCO images & question-type answer prior (text side) & free-form answer & 8,000 questions from the VQA-CP train split & 1,000 (same split, in-distribution) & 2,000 in-distribution + 3,000 from the VQA-CP test split (shifted priors); image-disjoint \\
Alignment captions & MS-COCO captions & -- & -- & 10,000 pairs & -- & -- \\
\end{longtable}
}

\textbf{Sampling.}
- \emph{Waterbirds:} official splits, kept unchanged for comparability with the group-robustness literature
\cite{sagawa2020groupdro,kirichenko2023dfr}.
- \emph{Watermark:}
- Training uses stratified sampling with 95\% of examples shortcut-aligned (watermark ↔ dog), the same
strength as Waterbirds.
- Validation and test are \textbf{group-balanced}, so every group is measured with equal precision.
- All splits are deterministic (\texttt{study/build\_tasks.py}, seed 0).
- \emph{CelebA:} train and validation are stratified subsamples that preserve group proportions, to fit the
compute budget; the official test set is used in full.

\textbf{Inclusion and exclusion.}
- Alignment captions exclude every image in any task's validation or test split, to prevent leakage.
- The Watermark label comes from the COCO object annotations: an image is positive if any annotation is
``dog''.
- No run was excluded after the fact. Runs that crashed (out-of-memory, session limits) were re-run with the
same seed, and every finished run is reported (\texttt{results\_study/*/all\_runs.csv}).

\textbf{Power analysis (sample size = seeds).}
- The unit of replication is the training seed. Each seed varies data order, LoRA initialisation and
stage-2 sampling \cite{bouthillier2021variance}.
- Power for the one-sided paired t-test at α = 0.025 was computed from the non-central t distribution
\cite{cohen1988power}.

{\def\LTcaptype{none} % do not increment counter
\begin{longtable}[]{@{}
  >{\centering\arraybackslash}p{(\linewidth - 10\tabcolsep) * \real{0.1667}}
  >{\centering\arraybackslash}p{(\linewidth - 10\tabcolsep) * \real{0.1667}}
  >{\centering\arraybackslash}p{(\linewidth - 10\tabcolsep) * \real{0.1667}}
  >{\centering\arraybackslash}p{(\linewidth - 10\tabcolsep) * \real{0.1667}}
  >{\centering\arraybackslash}p{(\linewidth - 10\tabcolsep) * \real{0.1667}}
  >{\centering\arraybackslash}p{(\linewidth - 10\tabcolsep) * \real{0.1667}}@{}}
\toprule\noalign{}
\begin{minipage}[b]{\linewidth}\centering
Paired effect size d\_z
\end{minipage} & \begin{minipage}[b]{\linewidth}\centering
Power with 3 seeds
\end{minipage} & \begin{minipage}[b]{\linewidth}\centering
5 seeds
\end{minipage} & \begin{minipage}[b]{\linewidth}\centering
10 seeds
\end{minipage} & \begin{minipage}[b]{\linewidth}\centering
15 seeds
\end{minipage} & \begin{minipage}[b]{\linewidth}\centering
Seeds for 80\% power
\end{minipage} \\
\midrule\noalign{}
\endhead
\bottomrule\noalign{}
\endlastfoot
0.8 (large) & 0.13 & 0.28 & 0.62 & 0.82 & 15 \\
1.0 & 0.18 & 0.40 & 0.80 & 0.95 & 10 \\
1.5 & 0.32 & 0.71 & 0.99 & 1.00 & 6 \\
2.0 & 0.47 & 0.91 & 1.00 & 1.00 & 5 \\
\end{longtable}
}

The initial 3-seed design had power below 0.5 even for large effects. The gate comparison was therefore
extended to \textbf{15 seeds per arm on both tasks}, enough for 80\% power at d\_z ≥ 0.8. Observed effects
(Section 4 of the results) are d\_z = 1.10 and 1.34, giving achieved power of 0.98 and \textgreater{} 0.99. Phi-4-mini
Watermark (d\_z ≈ 1.0, power 0.40 at 5 seeds) was extended to 10 seeds \emph{after} its 5-seed test missed
α = 0.025 (p = 0.044). This data-dependent extension is disclosed, and the new seeds 5--9 are also tested on
their own as fresh data (p = 0.023).

\subsection{Variables and instrumentation}

\textbf{Independent variable: connector}, four levels:

{\def\LTcaptype{none} % do not increment counter
\begin{longtable}[]{@{}
  >{\raggedright\arraybackslash}p{(\linewidth - 6\tabcolsep) * \real{0.2222}}
  >{\raggedright\arraybackslash}p{(\linewidth - 6\tabcolsep) * \real{0.2222}}
  >{\centering\arraybackslash}p{(\linewidth - 6\tabcolsep) * \real{0.2778}}
  >{\centering\arraybackslash}p{(\linewidth - 6\tabcolsep) * \real{0.2778}}@{}}
\toprule\noalign{}
\begin{minipage}[b]{\linewidth}\raggedright
Level
\end{minipage} & \begin{minipage}[b]{\linewidth}\raggedright
Design
\end{minipage} & \begin{minipage}[b]{\linewidth}\centering
Image tokens
\end{minipage} & \begin{minipage}[b]{\linewidth}\centering
Trainable params
\end{minipage} \\
\midrule\noalign{}
\endhead
\bottomrule\noalign{}
\endlastfoot
MLP projector \cite{liu2024llava15} & LayerNorm--Linear--GELU--Linear over 2×2-pooled patches & 49 & 8.5M \\
Q-Former \cite{li2023blip2} & 32 learned queries, 8 heads, 1 block & 32 & 11.4M \\
Cross-attention, learned gate \cite{alayrac2022flamingo} & prompt tokens attend to patches; \texttt{e\ ←\ e\ +\ tanh(α)·attn}, α₀ = 0 & 0 & 8.9M \\
Cross-attention, fixed gate & identical block; gate fixed open, output layer zero-initialised & 0 & 8.9M \\
\end{longtable}
}

The two cross-attention levels have the \textbf{same parameters and FLOPs}; only whether \texttt{α} is trainable
differs. Both start as the identity map. This isolates the gate as the manipulated mechanism and removes
capacity as a rival explanation.

\textbf{Dependent variables.}
- \textbf{Primary: test worst-group accuracy (WGA)} \cite{sagawa2020groupdro}.
- Mean over three question phrasings, one of them never seen in training.
- Checkpoint chosen by validation WGA.
- \textbf{Secondary:}
- average accuracy: the plain mean over test images. It is \emph{not} reweighted to the training-group
frequencies, as some Waterbirds papers do \cite{sagawa2020groupdro}, so it is not directly comparable to their
``average accuracy'';
- ID and OOD accuracy and their gap \cite{koh2021wilds};
- balanced accuracy;
- seed variance of WGA (H2);
- WGA with the label-free checkpoint rule (validation average accuracy).
- \textbf{Mechanism measures:} the learned gate value \texttt{tanh(α)} and its gradient, logged every 10 alignment steps.
Up to round 4 the logged gradient is the value after global-norm clipping at 1 (as in Flamingo); from round 5
the raw pre-clip gradient and the pre-clip global norm are logged as well (\texttt{grad\_raw}, \texttt{grad\_norm\_total}).
- \textbf{Behavioural pathway test:} from round 5, every selected checkpoint is re-tested with the connector's visual
contribution removed (\texttt{ablated\_mean\_over\_templates}). If WGA does not change, the pathway is behaviourally
closed, the criterion of \cite{farazi2026gate}.
Also the share of attention on the watermark corner.

\textbf{Controlled variables, and why each is controlled.}
- \textbf{LLM:} Qwen3.5-4B \cite{qwen2026qwen35}, NF4 4-bit with LoRA r = 16 \cite{dettmers2023qlora,hu2022lora}.
Controlled because LLM capacity and pretraining change shortcut reliance \cite{sagawa2020overparam}. Varied once
(Phi-4-mini \cite{abouelenin2025phi4mini}) to test generality.
- \textbf{Vision features:} frozen SigLIP \cite{zhai2023siglip}, cached once, so every connector receives
bit-identical inputs. Encoder choice strongly affects MLLMs \cite{tong2024cambrian,tong2024eyes}. Varied once
(SigLIP-large).
- \textbf{Data, prompts and answer read-out:}
- The same splits, the same chat template and the same 2 training phrasings for every connector.
- The prediction is the more likely of the two answer tokens, which removes free-text parsing as a
confounder.
- \textbf{Optimisation and budget:}
- \emph{Stage 1:} 300 alignment steps on 10k COCO captions with the LLM frozen.
- \emph{Stage 2:} 4 epochs, batch 16, learning rate 2e-4, cosine schedule.
- The same for every connector, with early stopping after 2 epochs without validation-WGA gain.
- The budget is equalised in steps, not wall-clock time, so faster connectors gain no extra training.
- \textbf{Image-token count:} differs by design (49 / 32 / 0). It is checked with a 36-token MLP ablation,
because token count is a known design factor \cite{cha2024honeybee,mckinzie2024mm1}.

\subsection{Procedure}

\begin{enumerate}
\def\labelenumi{\arabic{enumi}.}
\tightlist
\item
  \textbf{Data preparation} (\texttt{study/build\_tasks.py}, \texttt{study/cache\_features.py}): build splits and cache frozen
  patch features (one Kaggle CPU/GPU session per encoder).
\item
  \textbf{Stage 1, alignment} (\texttt{study/align.py}): train only the connector on COCO captions, with the LLM frozen.
  This follows the two-stage practice of \cite{liu2023llava,liu2024llava15}. 300 steps suffice because alignment
  is coarse and data-efficient \cite{groger2026rethinking}.
\item
  \textbf{Stage 2, task training} (\texttt{study/train\_task.py}): train LoRA and the connector on each task under ERM
  with natural (shortcut-biased) sampling.
\item
  \textbf{Stage 3, optional connector-only fixes:} freeze the LLM and LoRA, then retrain only the connector:

  \begin{itemize}
  \tightlist
  \item
    \textbf{Connector-DFR:} on a group-balanced half of validation, the other half selecting the model
    \cite{kirichenko2023dfr}.
  \item
    \textbf{Connector-JTT:} stage-2 errors upweighted ×20 \cite{liu2021jtt}.
  \end{itemize}
\item
  \textbf{Evaluation:}

  \begin{itemize}
  \tightlist
  \item
    Test metrics for each phrasing.
  \item
    Per-image predictions saved (\texttt{arrays.npz}) for the bootstrap.
  \item
    \texttt{environment.txt} (exact package versions and GPU) written next to every result.
  \end{itemize}
\item
  \textbf{Execution:}

  \begin{itemize}
  \tightlist
  \item
    Runs on Kaggle (2× NVIDIA T4, 16 GB), one job per GPU.
  \item
    Jobs are resumable; results are exported after every run.
  \item
    Seeds were assigned before the runs; no run was selected after seeing its result.
  \end{itemize}
\end{enumerate}

\textbf{Bias-minimisation steps.}
- Identical code paths for all connectors: one \texttt{train\_task.py} with a \texttt{-\/-connector} flag.
- A held-out question phrasing is used to detect prompt over-fitting.
- Two checkpoint-selection rules are reported: one uses group labels, one does not.
- Group-label use is stated for every method.
- The full per-seed results are released.

\subsection{Analysis strategy}

{\def\LTcaptype{none} % do not increment counter
\begin{longtable}[]{@{}
  >{\raggedright\arraybackslash}p{(\linewidth - 4\tabcolsep) * \real{0.3333}}
  >{\raggedright\arraybackslash}p{(\linewidth - 4\tabcolsep) * \real{0.3333}}
  >{\raggedright\arraybackslash}p{(\linewidth - 4\tabcolsep) * \real{0.3333}}@{}}
\toprule\noalign{}
\begin{minipage}[b]{\linewidth}\raggedright
Question
\end{minipage} & \begin{minipage}[b]{\linewidth}\raggedright
Test
\end{minipage} & \begin{minipage}[b]{\linewidth}\raggedright
Why this test
\end{minipage} \\
\midrule\noalign{}
\endhead
\bottomrule\noalign{}
\endlastfoot
H1: fixed \textgreater{} learned gate & One-sided \textbf{paired t-test} over seeds, with a \textbf{Wilcoxon signed-rank} test as the assumption-free check \cite{dror2018hitchhiker}; \textbf{Holm} correction across the four H1 cells (2 LLMs × 2 tasks); unpaired Welch and Mann--Whitney tests as checks & Seeds pair the two arms (same data order and initialisation), removing between-seed variance \\
H2: learned gate more variable & One-sided \textbf{variance-ratio F-test}, with the \textbf{Brown--Forsythe} test as a robust check, and \textbf{Fisher's exact test} on the number of collapsed runs (WGA \textless{} 50) & The F-test is sensitive to non-normality; Brown--Forsythe is not; the collapse count tests the practical claim directly \\
Pairwise connector differences & \textbf{Stratified bootstrap} of test images (2,000 resamples) and a \textbf{hierarchical bootstrap} over seeds and images \cite{efron1993bootstrap}; \textbf{Holm--Bonferroni} correction within each task \cite{holm1979simple} & The hierarchical bootstrap covers training variance, not just test-sample noise \\
``Connector type does not matter'' & \textbf{Two one-sided tests (TOST)}, margin ±3 points \cite{schuirmann1987tost} & A non-significant difference is not evidence of equivalence \\
Effect size & Paired \textbf{Cohen's d\_z} (mean difference / s.d. of differences), plus d\_av for comparison across designs \cite{cohen1988power} & p-values alone do not show practical importance \\
\end{longtable}
}

\textbf{Confirmatory vs exploratory.} Only the four H1 tests (2 LLMs × 2 tasks, Holm-corrected) and H2 on
Qwen3.5-4B are confirmatory. Every control (open and half-open gate, learning rate, alignment length, Group DRO,
injection ratio, SigLIP-large, CelebA, p-sweep, per-layer, fp16) is exploratory: its p-values are reported
without correction and are not used as evidence for the hypotheses.

\textbf{Hypothesis timeline (stated openly, since there is no public pre-registration).}
- 28 Oct 2026: 3-seed pilot of all connectors, including the fixed gate (commit \texttt{2c948b4}).
- 28 Oct 2026, after the pilot: seed-extension runs launched (Kaggle kernels \texttt{h1a}, \texttt{h1b}, \texttt{seeds34}).
- 29 Oct 2026: H1/H2 and their tests committed with the extended results (\texttt{7ab0563}, \texttt{study/hypotheses.py}).
- Because H1 was formed after seeing the pilot, it is also tested on the confirmation seeds only (seeds ≥ 3):
Waterbirds +10.5, t(11) = 4.06, p = 9.5 × 10⁻⁴, d\_z = 1.17; Watermark +26.7, t(11) = 4.86, p = 2.5 × 10⁻⁴,
d\_z = 1.40. The pilot alone was not significant (p = 0.16 and 0.12). (\texttt{results\_study/robustness\_checks.md})

\textbf{Assumption checks (reported in full).}
- \emph{Normality of the paired differences (Shapiro--Wilk):}
- Violated for Qwen Waterbirds (p = 0.009), which is why the Wilcoxon test is reported alongside;
it agrees (p = 3e-5).
- Not violated elsewhere (p ≥ 0.34).
- \emph{Normality of the learned-gate WGA:} violated on Qwen Waterbirds (p = 0.003), driven by the collapsed
seeds. The robust Brown--Forsythe test is therefore reported for H2.
- \emph{Independence:} each seed is a separate training run with its own random state. The two arms share only
the seed value, which is the intended pairing.

\subsection{Validity, reliability and ethics}

\textbf{Internal validity.}
- One factor varies at a time. The two gate arms share every parameter, but differ in two ways: whether \texttt{α}
trains, and where the zero initialisation sits (the learned gate starts with α = 0 and a default-initialised
\texttt{out\_proj}; the fixed gate starts with tanh α = 1 and \texttt{out\_proj} = 0).
- The second difference is forced, not chosen. If both α and \texttt{out\_proj} start at zero, the gradient of the
output with respect to α is proportional to \texttt{out\_proj} and that with respect to \texttt{out\_proj} to tanh α, so
both stay at zero and the visual pathway never trains. A learnable zero-initialised gate therefore needs a
non-zero \texttt{out\_proj}, as in Flamingo and LLaMA-Adapter.
- Two further arms separate trainability from the starting value. The design follows the residual-scaling
literature, where the start value of a learnable scale is the variable: ReZero and SkipInit start it at 0
\cite{bachlechner2021rezero,de2020skipinit}, and LayerScale at a small ε, reporting that ε trains better than 0
\cite{touvron2021cait}. Those scales are linear; Flamingo's gate passes through tanh, so its gradient is scaled
by 1 − tanh² α and vanishes as the gate saturates. Both arms are learnable, with \texttt{out\_proj} = 0 like the
fixed gate: \texttt{bottleneck\_open} starts at tanh α = 0.995 and \texttt{bottleneck\_half} at tanh α = 0.5. The open
gate's gradient is scaled by 1 − tanh² ≈ 0.01, so it is nearly fixed in practice; the half-open gate's is
scaled by 0.75, so it can move toward 0 or 1. The value 0.5 is our choice, not taken from prior work: it
is the midpoint of the gate's range at which the gradient is still within 25\% of its maximum (at tanh α = 0). If the half-open gate also matches the fixed gate, the harm is
tied to starting at zero; if it drifts shut and loses WGA, trainability itself is the problem.
- Pairing by seed shares the data order and the LoRA initialisation between the two arms of a seed (the
connector initialisations differ by design, see above). A seed does not replay a run exactly: GPU kernels
(4-bit matrix multiplies, linear-attention kernels, atomic additions) are non-deterministic, and library
versions drifted between jobs (transformers 5.21.0 to 5.22.1, recorded per job in \texttt{environment.txt}).
Rerunning seeds 0--1 changed WGA by 12.7 points on average for the learned gate and 5.2 for the fixed gate
(\texttt{results\_study/extra\_analyses.md} §9). Seeds are therefore treated as random replicates. Pairing is a
design choice, so unpaired Welch and Mann--Whitney tests are reported as well
(\texttt{results\_study/robustness\_checks.md}); they agree.
- The headline WGA averages three question phrasings, two of which were used in training. On the held-out
phrasing alone the gate effect is the same (+10.8, +25.5, +10.2 and +25.8; all p ≤ 0.006;
\texttt{extra\_analyses.md} §8).
- Model selection uses validation WGA (group labels), as is standard \cite{sagawa2020groupdro,kirichenko2023dfr};
the label-free rule (best validation average accuracy) is reported for every run and gives a larger gate
effect.
- Missing runs: Watermark seeds 8 and 14 were never launched because of the session budget; no finished run
was excluded. Both are being added so that each arm has 15 seeds.
- Rival explanations are tested directly: capacity (identical parameters), token count (36-token
ablation), encoder (SigLIP-large), LLM (Phi-4-mini) and training length (8 epochs).

\textbf{Construct validity.}
- Shortcut reliance is measured behaviourally, as WGA and the ID--OOD gap on test sets where the shortcut is
broken, not inferred from attention. Attention is not explanation \cite{jain2019attention}.
- The mechanism is measured directly: the logged gate value and its gradient and, from round 4, the injected
signal \texttt{‖tanh(α)·out\_proj(attn)‖\ /\ ‖e‖}, since a small gate can be offset by large output weights.
- Worst-group accuracy is reported with average and in-distribution accuracy, so a robustness gain cannot hide
an accuracy loss.

\textbf{External validity.}
- Coverage:
- three tasks: a natural background shortcut, a synthetic text shortcut and a demographic attribute;
- two LLM families;
- two encoder sizes;
- released 9B and 80B IDEFICS checkpoints, inspected for gate closure.
- Bounds:
- models are 4B scale or smaller and 4-bit;
- the vision encoder is frozen;
- the alignment budget is short.

\textbf{Reliability and replicability.}
- Deterministic splits.
- Exact scripts for every run in \href{../study/kaggle_kernels/}{\texttt{study/kaggle\_kernels/}}.
- Requirements (\href{../requirements-study.txt}{\texttt{requirements-study.txt}}, Kaggle pins in
\href{../study/requirements-kaggle.txt}{\texttt{study/requirements-kaggle.txt}}) and an \texttt{environment.txt} per run.
- Not bit-reproducible: GPU kernels are non-deterministic and Kaggle jobs installed the latest
\texttt{transformers}; same-seed reruns differ (§6, internal validity), so seeds are random replicates.
- 23 unit tests.
- All per-seed results and figure scripts.

\textbf{Ethics.}
- All data are public research datasets used under their terms:
- Waterbirds \cite{sagawa2020groupdro}: CUB-200-2011 birds \cite{wah2011cub} on Places backgrounds \cite{zhou2017places};
both are for non-commercial research.
- MS-COCO 2017 \cite{lin2014coco}: annotations CC BY 4.0; images under Flickr terms of use.
- CelebA \cite{liu2015celeba}: non-commercial research purposes only.
- VQA-CP v2 \cite{agrawal2018vqacp}: re-split of the VQA v2 annotations (CC BY 4.0) on COCO images. The JSON files
were taken from a public Kaggle mirror because the original host was unreachable; the mirror was checked
against the published split sizes (438,183 train and 219,928 test questions).
- Models: Qwen3.5-4B (Apache-2.0), Phi-4-mini-instruct (MIT), SigLIP (Apache-2.0). Check each model card
for the exact version used (\texttt{environment.txt}).
- SigLIP was pre-trained on web data (WebLI) that may contain CUB, Places or COCO images; this affects every
connector equally, because the encoder is frozen and shared.
- No new human-subject data were collected, so no IRB approval was required.
- CelebA contains face images and a binary gender attribute. It is used only as the standard
group-robustness benchmark \cite{sagawa2020groupdro}; no identity inference is performed.
- Compute: about 150 T4 GPU-hours of task training across 222 logged runs (from \texttt{train\_minutes}), plus data
preparation, alignment and evaluation, on Kaggle free-tier accounts. Rounds 3--5 (controls and extensions)
add roughly 120--150 GPU-hours. Preliminary work (the VLM-SHIELD pilot, failed or cancelled sessions) is not
in these totals and should be estimated by the authors for the NeurIPS-style compute statement.
